% !TEX program = pdflatex
\documentclass[11pt]{article}

% --- Preamble ---
\usepackage[margin=1in]{geometry}
\usepackage{fancyhdr}
\usepackage{hyperref}
\usepackage{xcolor}
\usepackage{graphicx}
\usepackage{booktabs}
\usepackage{amsmath}
\usepackage{amssymb}
\usepackage{amsthm}
\usepackage{enumitem}
\usepackage{listings}
\usepackage{caption}
\usepackage{subcaption}
\usepackage{float}
\usepackage{textcomp}
\usepackage{titling}
\usepackage{todonotes}

% --- Colors ---
\definecolor{linkblue}{RGB}{0, 102, 204}
\definecolor{codebg}{RGB}{245, 245, 245}
\definecolor{codeframe}{RGB}{220, 220, 220}

% --- Convenience commands ---
\newcommand{\menu}[1]{\textsf{#1}}

% --- Hyperref setup ---
\hypersetup{
    colorlinks=true,
    linkcolor=linkblue,
    urlcolor=linkblue,
    citecolor=linkblue,
    pdftitle={NeoSetzer Example Project},
    pdfauthor={NeoSetzer Contributors},
}

% --- Header / Footer ---
\pagestyle{fancy}
\fancyhf{}
\fancyhead[L]{\small\slshape NeoSetzer Example Project}
\fancyhead[R]{\small\thepage}
\renewcommand{\headrulewidth}{0.4pt}
\renewcommand{\footrulewidth}{0pt}

% --- Theorem environments ---
\theoremstyle{definition}
\newtheorem{definition}{Definition}[section]
\newtheorem{example}[definition]{Example}
\newtheorem{remark}[definition]{Remark}

\theoremstyle{plain}
\newtheorem{theorem}[definition]{Theorem}
\newtheorem{lemma}[definition]{Lemma}
\newtheorem{proposition}[definition]{Proposition}

% --- Listings setup ---
\lstset{
    basicstyle=\ttfamily,
    backgroundcolor=\color{codebg},
    frame=single,
    framerule=0.6pt,
    rulecolor=\color{codeframe},
    numbers=left,
    numberstyle=\tiny\color{gray},
    stepnumber=1,
    breaklines=true,
    captionpos=b,
    xleftmargin=1.5em,
    xrightmargin=1.5em,
    showspaces=false,
    showstringspaces=false,
    showtabs=false,
}

% The `listings` package does not ship a JSON language definition.
% Define a minimal JSON language so the sample configuration highlights correctly.
\lstdefinelanguage{JSON}{
    morekeywords={true,false,null},
    sensitive=false,
    morecomment=[l]{//},
    morestring=[b]",
    morestring=[b]',
    alsoletter={:},
}

% --- Title ---
% Push the title block toward the vertical middle of the cover page.
\pretitle{\begin{center}\vspace*{0.25\textheight}\Huge\bfseries}
\posttitle{\end{center}}
\predate{\begin{center}\large}
\postdate{\par\vspace*{0.2\textheight}\end{center}}

\title{NeoSetzer Example Project: A Tour of \LaTeX{} and Editor Features}
\author{NeoSetzer Contributors}
\date{\today}

% --- Document ---
\begin{document}

\maketitle
\thispagestyle{empty}

\clearpage

\tableofcontents
\newpage

% --- Abstract ---
\begin{abstract}
This document is a compact showcase of common \LaTeX{} elements
together with NeoSetzer-specific features and tips.
It is intended as a friendly starting point: you can read the source
to see how each effect is produced, and modify it freely.
Sections cover formatting, mathematics, lists, tables, figures,
code listings, cross-references, and practical guidance for using
the NeoSetzer editor.
\end{abstract}

% ============================================================
\section{Introduction}
\label{sec:intro}

Welcome to \LaTeX{}!
This guide walks through the most useful building blocks of a typical
document.
Each subsection is deliberately short so you can jump straight to the
part you need.

\subsection{What this document contains}
\label{sec:intro:contains}

As summarized in Table~\ref{tab:fruits}, this example covers:

\begin{itemize}[noitemsep]
    \item Text formatting and structure
    \item Mathematics, from inline formulas to multi-line alignments
    \item Various list styles
    \item Professional tables and figures
    \item Code listings
    \item Cross-references, footnotes, and hyperlinks
    \item NeoSetzer editor features and tips
\end{itemize}

% The following chapters are separate files on purpose. Open one in
% NeoSetzer to see the root-document workflow in action.
% !TEX root = ../main.tex

\section{Getting Started with NeoSetzer}
\label{sec:getting-started}

This folder is your own copy of the NeoSetzer example project. It is safe to
edit, save, build, move, or delete: later NeoSetzer updates do not overwrite
it. Build \texttt{main.tex} first, then keep this document open while you
explore the source and the PDF preview.

\subsection{Choose a starting point}
\label{sec:getting-started:starting-point}

The welcome screen offers several paths into a document. Create a blank LaTeX
file when you want a minimal starting point, or use the \textbf{Document
Wizard} to choose a document class, packages, title information, and other
sensible defaults. The wizard also keeps its small configuration presets
separate from \textbf{user document templates}: a user document template is a
snapshot of a complete LaTeX source document that can be saved, previewed, and
used as a starting point later.

This example project is deliberately another starting point. It is useful when
you prefer to learn by changing a working document rather than assembling every
piece from scratch.

\subsection{Build a first PDF}
\label{sec:getting-started:first-build}

The Magic Comment at the beginning of \texttt{main.tex} selects
\texttt{pdflatex}. NeoSetzer reads this small source-level instruction and
uses it only when the configured program is supported. The document keeps
PDFLaTeX as a portable baseline; you can change the configured build system in
\menu{Preferences} when your project needs XeLaTeX, LuaLaTeX, Tectonic, or
another available builder.

Use \textbf{Save and Build} after changing one sentence. The default shortcut
is commonly F5, but keyboard shortcuts are configurable, so the authoritative
list is \menu{Preferences \textrightarrow{} Keyboard Shortcuts}. The preview
shows the generated PDF, and the build log identifies errors, warnings, and bad
boxes with links back to source locations.

\subsection{A small task to try}
\label{sec:getting-started:task}

Change the title in \texttt{main.tex}, save, and build again. Then open the
Document Structure sidebar and select this section. It should take you to the
corresponding source heading. This short loop---edit, build, preview, and
navigate---is the foundation for the rest of the project.

% !TEX root = ../main.tex

\section{Writing, Structure, and Navigation}
\label{sec:writing-navigation}

NeoSetzer continuously reads the source so that the editor can offer
completion, navigation, and structure information while you write. These
features are source-level helpers: they are designed to be fast and reliable,
not to emulate every possible TeX macro expansion.

\subsection{A \textit{nested} heading title}
\label{sec:writing-navigation:nested-title}

This heading intentionally contains a nested command. It is a useful sample
for the Document Structure sidebar: the title remains visible as one heading,
while its italic formatting belongs to the rendered PDF. Try selecting it from
the sidebar and compare the source, outline, and preview.

\subsection{Labels, references, and completion}
\label{sec:writing-navigation:labels}

This section has a label, so another part of the document can refer to
Section~\ref{sec:writing-navigation}. NeoSetzer indexes labels, references,
citations, and bibliography keys from open documents to improve completion.
Type a backslash for command completion, or use the configurable manual
completion shortcut when you prefer to invoke the list explicitly.

\subsection{TODO items and comments}
\label{sec:writing-navigation:todos}

The following visible note is also a small navigation sample:
\todo{Replace this note after exploring the TODO list in Document Structure.}

NeoSetzer recognizes \texttt{\textbackslash todo\{...\}} items and lists them
in the document-structure sidebar. Select the TODO entry to jump to its source
location. In a real project you can delete the command after finishing the
task; this example keeps it as a discoverable demonstration.

\subsection{Open a child file directly}
\label{sec:writing-navigation:root}

The first line of this file contains \texttt{\% !TEX root = ../main.tex}.
Open this child file in NeoSetzer and build: the relative root declaration
points the build at \texttt{main.tex}, while keeping the child convenient to
edit. This pattern is useful for reports, theses, and books split into chapters.

% !TEX root = ../main.tex

\section{Project Workflows}
\label{sec:project-workflows}

A larger LaTeX project is more than one source file. NeoSetzer tracks common
project dependencies such as \texttt{\textbackslash input},
\texttt{\textbackslash include}, bibliography files, local letter options, and
custom document classes so that related files remain discoverable from the
workspace.

\subsection{Root documents and Magic Comments}
\label{sec:project-workflows:magic-comments}

The root declaration at the top of this child file is intentionally literal:
\begin{lstlisting}[language={[LaTeX]TeX}, caption={A child-file build declaration.}]
% !TEX root = ../main.tex
\end{lstlisting}

Magic Comments are metadata, not shell commands. NeoSetzer accepts a small,
safe subset: supported \texttt{program} values select a known LaTeX engine,
and a relative \texttt{root} value identifies a \texttt{.tex} document to
build. The application does not execute arbitrary source comments.

\subsection{Bibliography files}
\label{sec:project-workflows:bibliography}

The citation in the main document uses the bundled \texttt{references.bib}
file. Build the project twice, or let your configured build tool manage the
passes, so BibTeX can create the bibliography. The result demonstrates a real
multi-file dependency rather than a citation only written as prose.

\subsection{Structure numbering}
\label{sec:project-workflows:numbering}

The Document Structure sidebar displays standard source-level section numbers
when they can be determined reliably. It observes starred section commands and
literal section-counter changes, but it intentionally does not pretend to run
arbitrary TeX macros or custom number-format redefinitions.

The appendix at the end of \texttt{main.tex} demonstrates the common
\texttt{\textbackslash appendix} transition. In an article, the first appendix
section appears as \texttt{A}; its child headings appear as \texttt{A.1}.

\subsection{Build diagnostics}
\label{sec:project-workflows:diagnostics}

When a build fails, select an item in the build log to jump back to the
reported source line. For a harmless experiment, add an unknown command on a
new line, build once, inspect the diagnostic, and then undo the edit. Keep
experiments in this user-owned copy rather than in a valuable production
project.

% !TEX root = ../main.tex

\section{Tables and Imported Data}
\label{sec:tables-data}

The earlier tables show ordinary LaTeX source. NeoSetzer also provides an
\textbf{Insert Table} dialog for creating a grid, selecting alignment and
borders, merging constrained cells, and editing cell contents before inserting
LaTeX into the document.

\subsection{CSV and TSV practice data}
\label{sec:tables-data:csv}

This project includes \texttt{data/example-table.csv}. Open it in a text
editor, copy its rows, then open NeoSetzer's \menu{Insert Table} dialog and
choose \menu{Paste TSV/CSV}. The dialog detects tab-separated text first and
otherwise treats comma-separated CSV as the default; it also provides explicit
CSV delimiter choices and an \menu{Import CSV/TSV File} action.

The CSV file is deliberately a \emph{practice input} for the table generator.
It is not an instruction for LaTeX to read CSV automatically. After the dialog
inserts the resulting tabular code, build this project to see the table in the
PDF and adjust the generated source as you would in your own document.

\subsection{Long tables and cell structure}
\label{sec:tables-data:longtable}

For a table that may span pages, choose the \texttt{longtable} environment in
the generator or adapt the generated source. The existing table examples use
\texttt{booktabs}; they demonstrate professional horizontal rules and a
\texttt{\textbackslash multicolumn} total row. The editor's table workflow is
meant to speed up correct starting code, not to hide the LaTeX source from you.

\subsection{A suggested experiment}
\label{sec:tables-data:experiment}

Import \texttt{data/example-table.csv}, insert the result below this paragraph,
and add a caption and label. Then refer to it from a later paragraph with
\texttt{\textbackslash ref}. This links the table generator, labels,
completion, PDF preview, and Document Structure into one small reproducible
exercise.


% ============================================================
\section{Text Formatting}
\label{sec:text}

\LaTeX{} provides several levels of emphasis.
You can write \emph{emphasized} text, \textbf{bold} text,
\textit{italic} text, \textsc{small caps}, and \underline{underlined}
text.
Combinations such as \textbf{\emph{bold italic}} also work.

\subsection{Lists}

\subsubsection{Itemize}
\label{sec:text:lists:itemize}

A simple bullet list:

\begin{itemize}[noitemsep]
    \item First item
    \item Second item
    \item Third item with a longer description that spans multiple lines
          and demonstrates how itemize wraps text naturally
\end{itemize}

\subsubsection{Enumerate}
\label{sec:text:lists:enumerate}

A numbered list:

\begin{enumerate}[noitemsep]
    \item Step one
    \item Step two
    \item Step three
\end{enumerate}

\subsubsection{Description}
\label{sec:text:lists:description}

A description list is handy for defining terms:

\begin{description}[noitemsep]
    \item[NeoSetzer] A \LaTeX{} editor built with GTK and libadwaita.
    \item[\LaTeX] A high-quality typesetting system for technical documents.
    \item[PDF] Portable Document Format, the usual output of \LaTeX{}.
\end{description}

\subsection{Quotes and Remarks}
\label{sec:text:quotes}

Use \texttt{quote} for short quotations and \texttt{quotation} for
longer ones.
Here is a short quote:

\begin{quote}
    \small
    The \TeX{} typesetting system was designed by Donald Knuth
    and first released in 1978.
\end{quote}

A remark environment is useful for side notes:

\begin{remark}
    You do not need to memorize all \LaTeX{} commands.
    Cheat sheets and editors with autocompletion help a lot.
\end{remark}

% ============================================================
\section{Mathematics}
\label{sec:math}

\LaTeX{} excels at typesetting mathematics.
This section shows inline formulas, displayed equations, and
multi-line alignments.

\subsection{Inline Math}
\label{sec:math:inline}

The quadratic formula is $x = \frac{-b \pm \sqrt{b^2 - 4ac}}{2a}$.
You can also use \(\frac{\partial u}{\partial t}\) inline without
breaking the paragraph flow.

\subsection{Displayed Equations}
\label{sec:math:display}

Einstein's famous relation:

\begin{equation}
    \label{eq:einstein}
    E = mc^2
\end{equation}

A system of linear equations:

\begin{align}
    \label{eq:linear}
    f(x) &= (x+a)(x+b) \nonumber \\
         &= x^2 + (a+b)x + ab
\end{align}

A piecewise definition:

\begin{equation}
    \label{eq:piecewise}
    |x| =
    \begin{cases}
        x,  & \text{if } x \ge 0, \\
        -x, & \text{if } x < 0.
    \end{cases}
\end{equation}

A matrix:

\begin{equation}
    \label{eq:matrix}
    A =
    \begin{pmatrix}
        a_{11} & a_{12} & \cdots & a_{1n} \\
        a_{21} & a_{22} & \cdots & a_{2n} \\
        \vdots & \vdots & \ddots & \vdots \\
        a_{m1} & a_{m2} & \cdots & a_{mn}
    \end{pmatrix}
\end{equation}

\subsection{Theorems and Proofs}
\label{sec:math:theorems}

\begin{theorem}[Pythagoras]
    \label{thm:pythagoras}
    In a right-angled triangle, $a^2 + b^2 = c^2$, where $c$ is the
    hypotenuse.
\end{theorem}

\begin{proof}
    This is the classic Euclidean proof.
    Many algebraic proofs also exist.
\end{proof}

\begin{lemma}
    \label{lem:trichotomy}
    For any real number $x$, exactly one of the following holds:
    $x > 0$, $x = 0$, or $x < 0$.
\end{lemma}

\begin{definition}[Prime Number]
    \label{def:prime}
    A prime number is a natural number greater than $1$ that has no
    positive divisors other than $1$ and itself.
\end{definition}

\begin{example}
    \label{ex:primes}
    The first six prime numbers are $2, 3, 5, 7, 11, 13$.
\end{example}

% ============================================================
\section{Tables}
\label{sec:tables}

Table~\ref{tab:fruits} uses the \texttt{booktabs} package for
professional horizontal rules.

\begin{table}[H]
    \centering
    \caption{Example Table}
    \label{tab:fruits}
    \begin{tabular}{llr}
        \toprule
        \textbf{Item} & \textbf{Description} & \textbf{Price} \\
        \midrule
        Apple  & Fresh fruit      & \$1.00 \\
        Banana & Yellow fruit     & \$0.50 \\
        Cherry & Small red fruit  & \$2.00 \\
        \midrule
        \multicolumn{2}{r}{\textbf{Total:}} & \$3.50 \\
        \bottomrule
    \end{tabular}
\end{table}

A second table with aligned columns:

\begin{table}[H]
    \centering
    \caption{Notation Summary}
    \label{tab:notation}
    \begin{tabular}{p{3.5cm} p{6cm}}
        \toprule
        \textbf{Symbol} & \textbf{Meaning} \\
        \midrule
        $\mathbb{R}$  & The set of real numbers \\
        $\forall$     & Universal quantifier, ``for all'' \\
        $\exists$     & Existential quantifier, ``there exists'' \\
        $\implies$    & Logical implication \\
        $\iff$        & Logical equivalence \\
        $\partial$    & Partial derivative \\
        $\nabla$      & Gradient operator \\
        $\infty$      & Infinity \\
        $\emptyset$   & Empty set \\
        $\approx$     & Approximately equal \\
        $\sim$        & Distributed as \\
        $\propto$     & Proportional to \\
        \bottomrule
    \end{tabular}
\end{table}

% ============================================================
\section{Figures and Captions}
\label{sec:figures}

Figure~\ref{fig:placeholder} shows a placeholder box.
In a real document, replace it with \texttt{\textbackslash
includegraphics}.

\begin{figure}[H]
    \centering
    \fbox{\parbox{0.55\textwidth}{\centering
        Figure placeholder\\
        Insert image with \texttt{\textbackslash includegraphics}
    }}
    \caption{Example Figure}
    \label{fig:placeholder}
\end{figure}

You can also place figures side by side using \texttt{subcaption}:

\begin{figure}[H]
    \centering
    \begin{subfigure}[b]{0.45\textwidth}
        \centering
        \fbox{\parbox{0.8\linewidth}{\centering Sub-figure A}}
        \caption{First sub-figure}
        \label{fig:sub:a}
    \end{subfigure}
    \hfill
    \begin{subfigure}[b]{0.45\textwidth}
        \centering
        \fbox{\parbox{0.8\linewidth}{\centering Sub-figure B}}
        \caption{Second sub-figure}
        \label{fig:sub:b}
    \end{subfigure}
    \caption{Multiple sub-figures}
    \label{fig:subfigures}
\end{figure}

% ============================================================
\section{Code Listings}
\label{sec:code}

The \texttt{listings} package renders source code with syntax-aware
highlighting.

\subsection{Python}
\label{sec:code:python}

\begin{lstlisting}[language=Python, caption={A simple Python function.}]
def fibonacci(n: int) -> int:
    """Return the n-th Fibonacci number."""
    if n < 0:
        raise ValueError("n must be non-negative")
    a, b = 0, 1
    for _ in range(n):
        a, b = b, a + b
    return a

if __name__ == "__main__":
    for i in range(10):
        print(f"F({i}) = {fibonacci(i)}")
\end{lstlisting}

\subsection{JSON}
\label{sec:code:json}

\begin{lstlisting}[language=JSON, caption={A sample configuration file.}]
{
    "name": "setzer-example",
    "version": "1.0.0",
    "dependencies": {
        "gtk": ">=4.0",
        "libadwaita": ">=1.0"
    },
    "settings": {
        "theme": "auto",
        "fontSize": 12,
        "autoSave": true
    }
}
\end{lstlisting}

% ============================================================
\section{Cross-References}
\label{sec:references}

This document is heavily cross-referenced so you can see how labels
and references work together.

\subsection{Links}
\label{sec:references:links}

Equation~(\ref{eq:einstein}) expresses mass--energy equivalence.
Theorem~\ref{thm:pythagoras} appears in Section~\ref{sec:math:theorems}.
Definition~\ref{def:prime} and Example~\ref{ex:primes} are in the same
section.
Table~\ref{tab:fruits} and Table~\ref{tab:notation} are in
Section~\ref{sec:tables}.
Figure~\ref{fig:placeholder} and the sub-figures
(\ref{fig:sub:a}, \ref{fig:sub:b}, \ref{fig:subfigures}) are in
Section~\ref{sec:figures}.

You can also create clickable hyperlinks:
\url{https://www.latex-project.org/} and
\href{https://github.com/}{GitHub}.

\subsection{Footnotes}
\label{sec:references:footnotes}

\LaTeX{} supports footnotes easily.\footnote{This is a footnote.}
You can add as many as you need, and \LaTeX{} will place them at the
bottom of the page.\footnote{Another footnote with a citation:
\cite{knuth1986}.}

% ============================================================
% !TEX root = ../main.tex

\section{NeoSetzer Feature Map}
\label{sec:feature-map}

This chapter is a map for readers who want to understand NeoSetzer beyond the
basic edit--build--preview loop. It is organized by task rather than by every
menu item. Settings, installed TeX tools, dictionaries, and keyboard shortcuts
can differ between systems, so use \menu{Preferences} and the Command Palette
as the live source of truth for your own installation.

\subsection{Find and run actions}
\label{sec:feature-map:actions}

Use the \textbf{Command Palette} when you know what you want to do but not
where the action lives. It groups commands by area, searches by name and
keywords, and explains when a command is unavailable in the current document
or view. This is often the fastest way to discover an operation before deciding
whether you want to assign or change a shortcut.

The welcome screen provides the other high-level entry points: create a blank
LaTeX or BibTeX document, open a file, start the Document Wizard, reopen recent
work, or create another writable example-project copy.

\subsection{Write efficiently}
\label{sec:feature-map:editing}

The editor combines syntax highlighting, automatic matching brackets,
auto-indentation, line operations, code folding, search and replace, and
context-aware completion. Completion adapts to LaTeX context: commands,
references, citations, labels, environments, symbols, enumerate items, and
Beamer frames are examples of suggestions that can be offered when the source
makes them meaningful. The \textbf{Shortcuts Bar} and \textbf{Symbols} sidebar
provide discoverable insertions when you prefer not to memorize commands.

Spell checking depends on an available dictionary and the language settings on
your system. It is best treated as a writing aid rather than a LaTeX validator:
use the build log for typesetting or package errors. Editor appearance, font,
line wrapping, indentation, current-line highlighting, and related behaviour
can be adjusted in the Appearance and Editor preference pages.

\subsection{Navigate, inspect, and organize}
\label{sec:feature-map:navigation}

The \textbf{Document Structure} sidebar lists headings, labels, project files,
and TODO entries. It shows ordinary section numbers when they are reliably
available from the source, including common appendices and literal counter
changes such as the sample in Appendix~\ref{sec:appendix-structure}. It does
not execute arbitrary TeX macros, so exotic number redefinitions may differ
from the final PDF.

Use the \textbf{Document Statistics} view when you need counts and reading
information, and use source labels and references to keep large documents
connected. The project and structure views are particularly useful with the
child files in this example project: opening a child does not prevent you from
building the root document.

\subsection{Build, diagnose, and review PDF output}
\label{sec:feature-map:build-preview}

NeoSetzer can use configured LaTeX build systems such as PDFLaTeX, XeLaTeX,
LuaLaTeX, Tectonic, and \texttt{latexmk} when they are installed. Save and
build is a common starting action, but shortcuts are configurable; inspect the
Keyboard Shortcuts preference page instead of relying on one fixed key.

The PDF preview supports zooming, searching within the PDF, and SyncTeX source
navigation when the selected engine produces the required synchronization data.
The build log groups diagnostics and lets you jump to reported source lines. If
a PDF changes outside NeoSetzer, the preview shows a persistent reload prompt
rather than silently leaving you with an old result.

\subsection{Keep projects safe and recoverable}
\label{sec:feature-map:recovery}

Auto-save can preserve recoverable editor state after a crash, while session
management and the recent-documents list help restore a working set. Recent
documents can be pinned. The Editor preferences also control how NeoSetzer
responds when source files change outside the application; enable automatic
reload only when that behaviour fits your workflow.

Auto-build can rebuild after editing stops, but it is deliberately a preference
rather than an assumption. It is most useful when the selected builder is fast
and the document is already in a healthy state.

\subsection{Tables, templates, and structured starting points}
\label{sec:feature-map:structured-tools}

The Document Wizard creates a new source document from configuration choices.
Its lightweight presets are separate from \textbf{user document templates},
which store complete LaTeX source snapshots for later preview, selection, and
reuse. The Insert Table dialog generates editable LaTeX grids, supports
constrained cell merges and longtable choices, and can import or paste UTF-8
CSV/TSV data. Chapter~\ref{sec:tables-data} contains a safe practice input.

\subsection{External tools and AI Fix}
\label{sec:feature-map:external-tools}

\textbf{AI Fix} is optional. When configured, it can hand a selected build
error and nearby source context to an external agent command. That agent runs
outside NeoSetzer and can modify files, so review its permissions, use trusted
project directories only, and keep ordinary backup or version-control habits.
Combining automatic external-file reload with auto-build can create a fast
workflow, but it also means external edits may be incorporated and rebuilt
quickly; enable that combination only when you understand the consequences.

\subsection{Where to tune behaviour}
\label{sec:feature-map:preferences}

\begin{description}[style=nextline]
    \item[Appearance]
        Select application and editor appearance, including themes and fonts.
    \item[Editor]
        Tune indentation, wrapping, matching behaviour, auto-save, external
        changes, and other writing-time behaviour.
    \item[Autocomplete]
        Choose completion behaviour and related suggestion settings.
    \item[Build System]
        Choose an installed builder and configure build-related behaviour such
        as automatic builds.
    \item[Keyboard Shortcuts]
        Inspect current bindings and change them to fit your workflow.
\end{description}

\subsection{A deliberate exploration path}
\label{sec:feature-map:explore}

For a complete first tour, use the Command Palette to find Document Structure,
Document Statistics, Symbols, Insert Table, Build Log, and Preferences. Then
return to this project, make one small change in each feature area, and build
again. The safest way to learn an editor is to connect each interface control
to a visible source or PDF result.


% ============================================================
\section{About NeoSetzer}
\label{sec:about}

NeoSetzer is maintained as a continuation fork of the Setzer project. This
example project follows NeoSetzer's current workflows and links to the
repositories where development and discussion take place.

\begin{itemize}[noitemsep]
    \item \textbf{NeoSetzer repository}: \url{https://github.com/Sam-Fic/NeoSetzer}
    \item \textbf{Original Setzer project}: \url{https://github.com/cvfosammmm/Setzer}
\end{itemize}

% ============================================================
\section{Conclusion}
\label{sec:conclusion}

This example demonstrates the core elements of a \LaTeX{} document and a
practical map of NeoSetzer. From here you can:

\begin{enumerate}[noitemsep]
    \item Open the source and change sections, labels, captions, or data.
    \item Use the feature map to find one editor or workspace tool to try.
    \item Build the project, inspect the PDF and structure sidebar, and keep
          the user-owned copy as a starting point for your own work.
\end{enumerate}

Happy \TeX ing!

% ============================================================
\appendix
% !TEX root = ../main.tex

\section{Appendix Structure Sample}
\label{sec:appendix-structure}

This appendix is part of the example project so you can inspect how the
Document Structure sidebar represents the usual article transition from numeric
sections to alphabetic appendix sections. The heading should be shown as
\texttt{A}, and the next subsection as \texttt{A.1} when source-level numbering
is available.

\subsection{A child appendix heading}
\label{sec:appendix-structure:child}

This subsection deliberately has no special formatting. It is a compact sample
of a normal child heading under an appendix section.

\setcounter{subsection}{2}
\subsection{A literal counter change}
\label{sec:appendix-structure:counter}

The literal \texttt{\textbackslash setcounter\{subsection\}\{2\}} command
above means this source heading is numbered \texttt{A.3}. NeoSetzer tracks this
small, explicit counter rule because it is deterministic in source text. More
complex macro-based numbering remains outside the live outline model.


% ============================================================
\bibliographystyle{plain}
\bibliography{references}

\end{document}
